\section{A joint observation-budget limit}
\label{sec:budget}

The preceding fixed-update identities do not permit substitution of an
increasing stopping count without further bounds. Here we give such bounds for a
specified last-iterate schedule. The argument uses the classical decreasing-step
framework of \citet{moulines2011nonasymptotic}, with constants controlled as the
fourth update moment grows. It is an existence result under known preprocessing,
with no optimality or novelty claim.

For each sufficiently large integer $n$, let $(W,J) \in \{0,1\}^2$ have genuine
nonnegative joint-cell probabilities, and write
\[
 \begin{aligned}
  p_n & = \E W \to p \in (0,1), & q_n     & = \E J = n^{-1/2}, \\
  a_n & = p_n(1 - p_n),         & b_n     & = q_n(1 - q_n),    \\
  c_n & = \Cov(W,J),            & c_n/b_n & \to \kappa.
 \end{aligned}
\]
Choose $a_0 > 0$ and $n_0$ such that $a_n \ge a_0$ and $q_n \le a_0/64$ for
$n \ge n_0$. Set $x = (W - p_n,J - q_n)^\top$ and
$S_n = \diag(\sqrt{a_n},\sqrt{b_n})$. The response is
$Y = \alpha + x^\top\beta^* + \varepsilon$, with fixed signal $\beta^*$,
$\E(\varepsilon \mid W,J) = 0$, and
$\E(\varepsilon^2 \mid W,J) = \sigma^2 \in (0,\infty)$. The centering, scales,
and intercept $\alpha$ are known and held fixed. Use exactly $n$ iid fresh
observations, zero initialization, and
\begin{equation}
 \label{eq:budget-schedule}
 B_n = 1,\qquad k_n
     = n,\qquad \tau_n
     = \lceil16/q_n\rceil,\qquad \eta_{n,j}
     = \frac{2}{\tau_n + j},\quad 0
     \le j
     < n.
\end{equation}
Thus $n$ is an actual processed-observation budget. All steps and the stopping
count are deterministic, and no iterate averaging is used.

\begin{proposition}[Rare-coefficient error at an observation budget]
 \label{prop:budget}
 Let $u_j = S_n(\beta_{n,j} - \beta^*)$, $m_j = \E u_j$, and decompose
 $\Cov(u_j) = D_j + N_j$, where $D_j$ is the covariance of the noiseless
 trajectory and $N_j$ is the covariance of propagated observation noise. Put
 $C_* = \|S_n\beta^*\|^2$. For every $n \ge n_0$ and $0 \le k \le n$,
 \begin{align}
  |\E\beta_{n,R,k} - \beta_R^*|^2 & \le \frac{C_*}{b_n} \left( \frac{\tau_n}{\tau_n + k} \right)^{7/2},\label{eq:budget-bias} \\
  0 \le (D_k)_{RR}/b_n            & \le \frac{C_*}{b_n} \left( \frac{\tau_n}{\tau_n + k} \right)^3,\label{eq:budget-design}   \\
  0 \le (N_k)_{RR}/b_n            & \le \frac{4\sigma^2}{b_n(\tau_n + k)}.\label{eq:budget-noise}
 \end{align}
 Moreover, $nN_n \to (4\sigma^2/3)I_2$ in operator norm. Consequently,
 \begin{equation}
  \label{eq:budget-mse}
  nq_n\Var(\beta_{n,R,n}) \longrightarrow \frac{4\sigma^2}{3}, \qquad nq_n\E(\beta_{n,R,n} - \beta_R^*)^2
                          \longrightarrow \frac{4\sigma^2}{3}.
 \end{equation}
 The squared bias is $O(n^{-5/4})$, and the original-scale design variance is
 $O(n^{-1})$. These constants are uniform over bounded signals and the stated
 population class with fixed $a_0$ and $\sigma^2$.
\end{proposition}

\begin{proof}
 Suppress $n$ and put $z = S_n^{-1}x$, $H = \E zz^\top$, and
 $K = \E[\|z\|^2zz^\top]$. Binary feasibility gives $|c_n/b_n| \le 1$, so the
 off-diagonal entry of $H$ has magnitude at most $\sqrt{q_n/a_0} \le 1/8$.
 Pointwise, $z_C^2 \le 1/a_0$ and $z_R^2 \le 1/q_n$. Hence
 \[
  \tfrac78 I \preceq H
             \preceq \tfrac98 I,\qquad K
             \preceq \frac{65}{64q_n}H
             \preceq \frac2{q_n}I.
 \]
 For $s_j = \tau_n + j$, let $R_j = I - 2H/s_j$ and
 \[
  \mathcal V(U) = \E[(zz^\top-H)U(zz^\top-H)],\qquad \mathcal T_j(U)
                = R_jUR_j + \frac4{s_j^2}\mathcal V(U).
 \]
 Independence and conditional mean-zero noise give the exact recurrences
 \[
  \begin{split}
   m_{j+1} & = R_jm_j,\quad m_0 = -S_n\beta^*,                                          \\
   D_{j+1} & = \mathcal T_j(D_j) + \frac4{s_j^2}\mathcal V(m_jm_j^\top), \quad D_0 = 0, \\
   N_{j+1} & = \mathcal T_j(N_j) + \frac{4\sigma^2}{s_j^2}H,\quad N_0 = 0.
  \end{split}
 \]
 The two covariance contributions are positive semidefinite. In particular, the
 covariance $D_j$ is distinct from the noiseless second moment
 $U_j^{(0)} = m_jm_j^\top + D_j$, which propagates by $\mathcal T_j$.

 For $U \succeq 0$, the uniform bounds and $q_ns_j \ge 16$ imply
 \[
  \operatorname{tr}\mathcal T_j(U) \le \left( 1 - \frac7{2s_j} + \frac8{q_ns_j^2} \right)\operatorname{tr}U
                                   \le (1 - 3/s_j)\operatorname{tr}U.
 \]
 The product of these factors from $j = l$ to $k - 1$ is at most
 $\{(\tau_n + l)/(\tau_n + k)\}^3$ by an integral bound on the harmonic sum.
 This bounds $\operatorname{tr}U_k^{(0)}$ and therefore $D_k$. Separately,
 $\|R_j\| \le 1 - 7/(4s_j)$ gives the squared-mean exponent $7/2$. Finally,
 $\operatorname{tr}H = 2$ gives
 \[
  \operatorname{tr}N_{j+1} \le (1 - 3/s_j)\operatorname{tr}N_j + 8\sigma^2/s_j^2
                           \le 4\sigma^2/(s_j + 1)
 \]
 by induction from $N_0 = 0$. Dividing the rare-coordinate quantities by $b_n$
 proves the finite bounds.

 To evaluate the noise constant, let $\bar N_0 = 0$ and
 $\bar N_{j+1} = R_j\bar N_jR_j + 4\sigma^2H/s_j^2$. Then
 $F_j = N_j - \bar N_j \succeq 0$ and
 \[
  F_{j+1} = R_jF_jR_j + \frac4{s_j^2}\mathcal V(N_j).
 \]
 Since $\operatorname{tr}\mathcal V(U) \le (2/q_n)\operatorname{tr}U$ for
 $U \succeq 0$, each source has trace at most $32\sigma^2/(q_ns_j^3)$. The
 subsequent contraction product is exactly
 \[
  \prod_{l=j+1}^{k-1}(1 - 3/s_l) = \frac{(s_j - 2)(s_j - 1)s_j}
 {(\tau_n + k - 3)(\tau_n + k - 2)(\tau_n + k - 1)}.
 \]
 Thus
 \begin{equation}
  \label{eq:budget-amplification}
  0 \le \operatorname{tr}F_k \le \frac{32\sigma^2 k}{q_n(\tau_n + k - 3)^3}.
 \end{equation}
 In particular, $n\|F_n\| \to 0$ since $nq_n \to \infty$ and $\tau_n/n \to 0$.

 All $R_j$ commute with $H$, whose two eigenvalues $h$ tend to one. The
 corresponding eigenvalues of $\bar N_n$ are
 \[
  4\sigma^2h\sum_{j=0}^{n-1}s_j^{-2} \prod_{l=j+1}^{n-1}(1 - 2h/s_l)^2.
 \]
 Fix $\delta \in (0,1)$ and $m = \lceil\delta n\rceil$. The terms $j < m$ have
 total trace at most $4\sigma^2(\tau_n + m)^2/(\tau_n + n)^3$, by propagating
 $\bar N_m$ and using the preceding trace bound. After multiplication by $n$,
 their limsup is at most $4\sigma^2\delta^2$. For $m \le j < n$, the uniform
 logarithmic expansion is
 \[
  \log\prod_{l=j+1}^{n-1}(1 - 2h/s_l)^2 = 4h\log\frac{\tau_n + j + 1}{\tau_n + n}
                                          + O((\tau_n + j + 1)^{-1}).
 \]
 The product therefore converges uniformly to $(j/n)^4$ on this late interval.
 Multiplication by $n$ turns its sum into a Riemann sum with limit
 $4\sigma^2\int_\delta^1 x^2\,dx$. First taking $n \to \infty$, then
 $\delta \downarrow 0$, proves $n\bar N_n \to (4\sigma^2/3)I_2$.
 Equation~\eqref{eq:budget-amplification} gives the same limit for $N_n$.
 Finally, $b_n \sim n^{-1/2}$ and $\tau_n \sim 16n^{1/2}$ give the stated bias
 and design rates and \eqref{eq:budget-mse}.
\end{proof}

The result establishes the $(nq_n)^{-1/2}$ second-moment uncertainty scale, with
a schedule-specific constant. It gives neither a central limit theorem nor an
efficiency claim. The same bounds suffice when $q_n \to 0$, $nq_n \to \infty$,
and additionally $n^2q_n^3 \to \infty$; the last condition is sufficient for the
conservative design bound, not claimed necessary. For an auxiliary
iterated-limit comparison, extend the schedule to an infinite fresh-observation
stream. Taking rarity to zero at a fixed update count leaves the original-scale
mean at its zero initialization, while taking the update count to infinity at
fixed full-rank prevalence gives $\beta^*$. The joint path in the proposition
specifies which experiment is analyzed. It does not rank normalization choices
or imply a fixed-data sampling result.
